\begin{table}[H]
\centering
\caption{Standardized Effect Sizes for Main Outcomes}
\label{tab:sde}
\begin{adjustbox}{max width=\textwidth}
\begin{threeparttable}
\begin{tabular}{llcccccc}
\toprule
Outcome & Specification & $\hat{\beta}$ & SE & SD($Y$) & SDE & SE(SDE) & Classification \\
\midrule
\multicolumn{8}{l}{\textit{Panel A: Pooled}} \\
Log Employment & Border triple-diff & -0.0931 & 0.0208 & 1.917 & -0.0486 & 0.0109 & Small negative \\
Log Establishments & Border triple-diff & 0.0581 & 0.0199 & 1.562 & 0.0372 & 0.0127 & Small positive \\
Log Weekly Wage & Border triple-diff & -0.0379 & 0.0179 & 0.333 & -0.1138 & 0.0538 & Moderate negative \\
\midrule
\multicolumn{8}{l}{\textit{Panel B: Heterogeneous}} \\
Log Emp (Information) & Info vs exempt & -0.1319 & 0.0316 & 1.939 & -0.0680 & 0.0163 & Moderate negative \\
Log Emp (Professional) & Prof vs exempt & -0.0580 & 0.0153 & 1.855 & -0.0313 & 0.0083 & Small negative \\
\bottomrule
\end{tabular}
\begin{tablenotes}[flushleft]
\footnotesize
\item \textit{Notes:} \textbf{Country:} United States. \textbf{Research question:} Did the 2019 Illinois Supreme Court ruling in \textit{Rosenbach v.~Six Flags}, which eliminated the injury requirement for biometric privacy lawsuits under BIPA, reduce employment and business activity in biometric-exposed industries in Illinois relative to neighboring states? \textbf{Policy mechanism:} The ruling transformed BIPA from a dormant statute into the most actively litigated privacy law in the US by allowing any individual whose biometric data was collected without proper consent to sue for statutory damages (\$1,000--\$5,000 per violation), even without demonstrating actual harm; this created enormous litigation exposure for firms using fingerprint scanners, facial recognition, or other biometric technologies. \textbf{Outcome definition:} Log quarterly county-level employment (average of three monthly levels from QCEW) and log quarterly establishment counts, by NAICS sector. \textbf{Treatment:} Binary: Illinois counties in biometric-exposed industries (Information, Professional/Technical Services) versus BIPA-exempt industries (Finance under GLBA, Healthcare under HIPAA) after the January 25, 2019 ruling. \textbf{Data:} BLS Quarterly Census of Employment and Wages (QCEW), 2015Q1--2023Q4, county-sector-quarter panel for Illinois and five border states (Indiana, Wisconsin, Missouri, Iowa, Kentucky). \textbf{Method:} Triple-difference (Illinois $\times$ exposed industry $\times$ post-ruling) with county-sector and quarter fixed effects; standard errors clustered at the state level; wild cluster bootstrap for few-cluster inference. \textbf{Sample:} Border counties only (counties sharing a state boundary between Illinois and neighboring states) to sharpen geographic comparison. SDE $= \hat{\beta} / \text{SD}(Y)$ where SD($Y$) is the pre-treatment standard deviation. Classification refers to magnitude, not statistical significance: Large ($|$SDE$| > 0.15$), Moderate ($0.05$--$0.15$), Small ($0.005$--$0.05$), Null ($< 0.005$).
\end{tablenotes}
\end{threeparttable}
\end{adjustbox}
\end{table}

